% ============================================================================
% Abstract -- compact variant (~200 words, Nature-style single paragraph)
% ============================================================================

Masked language models pretrained on SMILES reduce de novo design to local,
token-level editing: mask part of a known ligand and resample it. We ask
whether the choice of \emph{which} tokens to mask carries structural
information, masking the atoms PLIP identifies as making protein--ligand
contacts in deposited co-crystal structures and comparing against uniformly
random masking of the same parents at the same rate. Because the pretrained
model infills such masks with low chemical fidelity, we fine-tune ChemBERTa
against a composite, non-differentiable objective over validity, drug-likeness,
synthetic accessibility, scaffold divergence, structural alerts and predicted
Tox21 activity, using two estimators as a controlled comparison: a REINFORCE
policy gradient with LoRA adapters and a KL anchor, and best-of-$K$ selection
followed by cross-entropy. Initial results are encouraging but preliminary ---
over three epochs, parseable generations rose from 20.9\% to 22.7\% and mean
composite reward from 0.101 to 0.110, with absolute validity still low. Part of
that ceiling appears to reflect the corpus rather than the objective: three
frameworks account for 28\% of distinct parents and crystallisation additives
for a quarter of all instances. We evaluate on a Bemis--Murcko framework
partition, and ongoing work targets generalisation across chemically diverse
molecular groups.
